% Printed 443.
\par\clearpage
\noindent\null\hfill{\small\textbf{HEIGHT MEASUREMENTS.}}\hfill\makebox[0pt][r]{\fontsize{13}{16}\selectfont 443}\par\medskip
\begin{center}
[3].\par\medskip
[The coefficient of refraction from the height measurements in the arc measurement\par
and its extension as far as Jever.]\par\bigskip
\textsc{Gauss} to \textsc{Schumacher}.\quad Göttingen, 27~December 1846.
\end{center}
\par\noindent\hspace*{\parindent}\textsc{Bessel} mentions [in a paper he left behind] my result for terrestrial refraction.
He thinks the reason I find a smaller coefficient than other astronomers is that I
observed only on sunny days, while they observed on both sunny and overcast days
without distinction. This statement is incorrect in several respects. First, the
factual claim does not hold without qualification. That paper was presumably written
before the book on the Prussian arc measurement. In the latter, p.~197, \textsc{Bessel} also
cites other results \textls[120]{between} which mine, as he gives it, falls. Moreover,
\textsc{Bessel} knew only the result of my measurements of 1821 and 1822, as printed in the
Jahrbuch for 1826. The later years give progressively larger coefficients, namely
\begin{center}
\begin{tabular}{r@{\ .\ .\ .\ }l}
1823&\(0{,}14125\)\\
1824&\(0{,}14778\)\\
1825&\(0{,}15826\).
\end{tabular}
\end{center}
\par\noindent\hspace*{\parindent}The mean result for all 5 years, from a number of sides of triangles whose total
length is 290 geographical miles, is \(0{,}13974\), which is greater than the result of
\textsc{Bessel}'s own observations. I can, however, readily explain that progression.
It is inaccurate to say that refraction is smaller on sunny days than on overcast
ones. The [reason] is instead that on sunny days refraction depends to a high degree
\textls[120]{on the time of day}, incomparably more than on overcast days. Refraction is
smallest at noon and during the morning and afternoon hours nearest noon. It increases
continuously with extraordinary regularity in the later afternoon hours, and towards
sunset it is certainly at least as large as on overcast days and more likely larger.
The time of least refraction is also when the air is least favourable for precise
\par\medskip\noindent\null\hfill\makebox[0.18\linewidth][l]{56\textsuperscript{*}}\par

% Printed 444.
\par\clearpage
\noindent\makebox[0pt][l]{\fontsize{13}{16}\selectfont 444}\hfill{\small\textbf{CORRESPONDENCE.\quad HEIGHT MEASUREMENTS.}}\hfill\null\par\medskip
\noindent observations. Since I regarded height measurements as a secondary matter, in the
first years I devoted to them chiefly the hours least favourable for the more precise
horizontal angle measurements, in particular the morning and midday hours. In the
following years I made less and less use of the morning hours and finally almost
stopped using them, so that I often measured zenith distances fairly late in the
afternoon as well. I attach very little weight to the reason \textsc{Bessel} gives for the
larger results obtained by the French astronomers, namely that \textsc{Borda} circles
give zenith distances that are too small. It could have only a very slight effect.
Moreover, my own measurements were also made with a \textsc{Borda} circle.
\begin{center}\rule{.16\linewidth}{.4pt}\end{center}
\begin{center}\textbf{REMARKS.}\end{center}
\begingroup\small
\par\noindent\hspace*{\parindent}With regard to the first part of the paper from \textsc{Bode}'s Astronomisches Jahrbuch,
the final values of the geographical positions have already been given on pp.~427/428.
The last of the two preceding letters to \textsc{Schumacher} was printed from the original,
and the first from a copy made from the original. Its numerical values could be
checked against the entries in one of \textsc{Gauss}'s handbooks. According to that
handbook, the values given for the Lichtenberg-Deister line in the table on p.~439
are to be replaced, in order, by \(1566{,}9\), \(222{,}3\), and \(0{,}1419\).
The height observations during the 5 years of the arc measurement yielded (according to
various entries in \textsc{Gauss}'s handbook) the following results for deriving the
coefficient of refraction:
\begin{center}
\setlength{\tabcolsep}{4pt}
\begin{tabular}{|l|r|r|r|r|}\hline
\multicolumn{1}{|c|}{\shortstack{Year of\\observations}}&\multicolumn{1}{c|}{\shortstack{\rule{0pt}{1pt}\\Number\\of lines}}&\multicolumn{1}{c|}{\shortstack{\rule{0pt}{1pt}\\Sum of the\\curvatures}}&\multicolumn{1}{c|}{\shortstack{\rule{0pt}{1pt}\\Sum of the\\refractions}}&\multicolumn{1}{c|}{\shortstack{\rule{0pt}{1pt}\\Coefficient of\\refraction}}\\\hline
1821 and 1822 (as far as Wulfsode-Wilsede)&28&\(33203^{\prime\prime}{,}272\)&\(4309^{\prime\prime}{,}986\)&\(0{,}12981\)\\
1823 (as far as Wilsede-Nindorf)&4&\(1778{,}089\)&\(178{,}572\)&\(0{,}12643\)\\
\phantom{1823 }(as far as Hamburg)&5&\(6105{,}219\)&\(862{,}415\)&\(0{,}14126\)\\
1824 (as far as Brillit-Garlste)&22&\(22549{,}844\)&\(3332{,}329\)&\(0{,}14778\)\\
1825 (as far as Varel-Jever)&8&\(7763{,}012\)&\(1228{,}541\)&\(0{,}15826\)\\\hline
\multicolumn{1}{|r|}{Total}&67&\(71399{,}436\)&\(9911{,}843\)&\(0{,}13882\)
\end{tabular}
\end{center}
\par\noindent\hspace*{\parindent}\textsc{Gauss} appears to have excluded the results of the first 4 lines of 1823 when
forming the mean on p.~443. In that case the coefficient of refraction is \(0{,}13980\).
\par\noindent\hfill\textsc{Krüger}.\par
\endgroup
\begin{center}\rule{.2\linewidth}{.4pt}\end{center}

% Printed 445.
\par\clearpage
\noindent\null\hfill{\small\textbf{POSTHUMOUS PAPERS.\quad HEIGHT MEASUREMENTS.}}\hfill\makebox[0pt][r]{\fontsize{13}{16}\selectfont 445}\par
\vspace*{2.4cm}
\begin{center}{\Large\textbf{\textls[140]{POSTHUMOUS PAPERS.}}}\par\vspace{1.3cm}
Terrestrial refraction.\par\medskip
[1.]
\end{center}
\par\noindent\hspace*{\parindent}It is assumed that the path of the light ray stays in a single plane and that
the directions of gravity at the various points of that path meet at a single
point \(C\). It is also assumed that the changes in direction along the path
are proportional to the changes in the angle at \(C\).
\par\noindent\hspace*{\parindent}Let \(P\) and \(P'\) be two [successive] points on the light ray's curve,
\(90^\circ+u\) the angle between \(PP'\) and \(PC\), and \(PC=r\). Let \(A\) be the
initial point of that line and \(B\) the endpoint. Let \(\theta\) be the angle
between \(CP\) and \(CA\), and \(n\mathrm{d}\theta\) the change in direction.
Finally, let \(a\) and \(b\) be the values of \(r\) at \(A\) and \(B\), and
\(\alpha\) and \(\beta\) the values of \(u\) at those points, i.e. \(\alpha\) the
elevation at which \(B\) appears from \(A\), and \(\beta\) the depression at
which \(A\) appears from \(B\).
\par\noindent\hspace*{\parindent}We have
\[u=\alpha+(1-n)\theta\]
\[\mathrm{d}r=r\mathrm{d}\theta.\operatorname{tang}u=r\mathrm{d}\theta.\operatorname{tang}(\alpha+(1-n)\theta)\]
\noindent [or]
\[\log r=-\tfrac{1}{1-n}\log\cos(\alpha+(1-n)\theta)+\mathrm{const.}\]
\[\tfrac{\mathrm{const.}}{r^{1-n}}=\cos(\alpha+(1-n)\theta).\]
\par\noindent\hspace*{\parindent}Consequently [since at \(r=a,\quad\theta=0\)]
\[\tfrac{a^{1-n}\cos\alpha}{r^{1-n}}=\cos(\alpha+(1-n)\theta).\]

% Printed 446.
\par\clearpage
\noindent\makebox[0pt][l]{\fontsize{13}{16}\selectfont 446}\hfill{\small\textbf{POSTHUMOUS PAPERS.}}\hfill\null\par\medskip
\begingroup
\setlength{\abovedisplayskip}{4pt}\setlength{\belowdisplayskip}{4pt}
\setlength{\abovedisplayshortskip}{2pt}\setlength{\belowdisplayshortskip}{2pt}
\par\noindent\hspace*{\parindent}If, then, at \(B\)
\[\theta=t\quad[\text{and }r=b],\]
\noindent then
\[\beta=\alpha+(1-n)t\]
\noindent [and]
\[\tfrac{a^{1-n}\cos\alpha}{b^{1-n}}=\cos\beta,\]
\noindent [so that]
\[\tfrac{\cos\alpha-\cos\beta}{\cos\alpha+\cos\beta}=\operatorname{tang}\tfrac12(\beta-\alpha)\operatorname{tang}\tfrac12(\beta+\alpha)=\tfrac{b^{1-n}-a^{1-n}}{b^{1-n}+a^{1-n}}.\]
\par\noindent\hspace*{\parindent}Set
\[\begin{aligned}\tfrac12(b+a)&=R+k\\b-a&=h,\end{aligned}\]
\noindent so
\[\begin{aligned}b&=R+k+\tfrac12h\\a&=R+k-\tfrac12h;\end{aligned}\]
\noindent [then:]
\[\operatorname{tang}\tfrac12(\beta-\alpha)\operatorname{tang}\tfrac12(\beta+\alpha)=\frac{\left(1+\tfrac{h}{2(R+k)}\right)^{1-n}-\left(1-\tfrac{h}{2(R+k)}\right)^{1-n}}{\left(1+\tfrac{h}{2(R+k)}\right)^{1-n}+\left(1-\tfrac{h}{2(R+k)}\right)^{1-n}},\]
\noindent or, to sufficient accuracy,
\[=\tfrac{(1-n)h}{2(R+k)}.\]
\par\noindent\hspace*{\parindent}Putting \(AB=\Delta\) for the length of the arc projected onto sea level gives
\[\Delta=Rt.\]
\par\noindent\hspace*{\parindent}Furthermore,
\[\text{I.}\hspace{3cm}\beta-\alpha=(1-n)t;\]
\noindent hence
\[\begin{aligned}
(1-n)h&=2(R+k)\operatorname{tang}\tfrac12(\beta-\alpha)\operatorname{tang}\tfrac12(\beta+\alpha)\\
&=\tfrac{\operatorname{tang}\frac12(\beta-\alpha)}{\frac12(\beta-\alpha)}.(R+k)(\beta-\alpha)\operatorname{tang}\tfrac12(\beta+\alpha),\\
h&=\tfrac{\operatorname{tang}\frac12(\beta-\alpha)}{\frac12(\beta-\alpha)}.\tfrac{R+k}{R}.\Delta.\operatorname{tang}\tfrac12(\beta+\alpha),
\end{aligned}\]
\noindent which may also be written
\[\text{II.}\hspace{3cm}h=\tfrac{R+k}{R}.\Delta.\operatorname{tang}\tfrac12(\beta+\alpha)\]
\par\noindent\hspace*{\parindent}When \(\alpha\) and \(\beta\) are both known, \(n\) follows from I and \(h\) from II.
\endgroup

% Printed 447.
\par\clearpage
\noindent\null\hfill{\small\textbf{HEIGHT MEASUREMENTS.}}\hfill\makebox[0pt][r]{\fontsize{13}{16}\selectfont 447}\par\medskip
\begingroup
\setlength{\abovedisplayskip}{5pt}\setlength{\belowdisplayskip}{5pt}
\begin{center}
[2.]\par\medskip
\begin{tabular}{l@{\ }l}
Depression of the horizon\ .\ .\ .\ .\ .\ .\ .&\(\delta\)\\
Radius of the Earth\ .\ .\ .\ .\ .\ .\ .\ .\ .\ .&\(R\)\\
Height of the eye\ .\ .\ .\ .\ .\ .\ .\ .\ .\ .\ .\ .\ .\ .&\(h\)\\
Amplitude of the terrestrial arc\ .\ .\ .\ .\ .\ .&\(v\)\\
Total refraction\ .\ .\ .\ .\ .\ .\ .\ .\ .\ .\ .\ .\ .&\(\varepsilon v\).
\end{tabular}
\end{center}
\par\noindent\hspace*{\parindent}[If we set]
\[R+h=r\]
\noindent [then:]
\[r^{1-\varepsilon}\cos(1-\varepsilon)v=R^{1-\varepsilon}.\]
\par\noindent\hspace*{\parindent}Hence [since \(\delta=(1-\varepsilon)v\), we have:]
\[\left(\tfrac{R+h}{R}\right)^{1-\varepsilon}.\cos\delta=1\]
\noindent [or]
\[\begin{aligned}
2\sin\tfrac12\delta^2&=1-\left(\tfrac{R}{R+h}\right)^{1-\varepsilon}\\
&=\tfrac{h}{R}(1-\varepsilon)-\tfrac{hh}{RR}\tfrac{(1-\varepsilon)(2-\varepsilon)}{2}\ .\ .\ .\ .
\end{aligned}\]
\par\noindent\hspace*{\parindent}If the path of the light is taken to be a circular arc, then
\[\begin{aligned}
\tfrac{R+h}{R}\left[=\tfrac{\sin(v-\delta)}{\sin v-\sin\delta}=\tfrac{\cos\frac12(v-\delta)}{\cos\frac12(v+\delta)}\right]\\
=\frac{\cos\tfrac{\varepsilon}{2-2\varepsilon}.\delta}{\cos\tfrac{2-\varepsilon}{2-2\varepsilon}.\delta}.
\end{aligned}\]
\begin{center}[3.]\end{center}
\par\noindent\hspace*{\parindent}\(A,\ B\)\ .\ .\ .\ .\ Heights of two points where the reciprocal depressions
\(a,\ b\) are observed.
\par\noindent\hspace*{\parindent}\(R\)\ .\ .\ .\ .\ Radius of curvature [of the Earth's sphere]; \(r\)\ .\ .\ .\ .\ Distance.
\begin{center}\(\mu\)\ .\ .\ .\ .\ Coefficient of refraction.\end{center}
\par\noindent\hspace*{\parindent}[Let \(v\) be the central angle corresponding to the distance \(r\), and \(\gamma\)
the central angle corresponding to the light ray. Then we have
\endgroup
